\documentclass{article} %210 mm × 297 mm
\usepackage[utf8]{inputenc}
\usepackage[a4paper,left=1.5cm, right=1.5cm, top=2cm, bottom=2cm]{geometry}
\usepackage{graphicx}
%\usepackage{blindtext}
\usepackage{multirow}
\usepackage{mhchem}
\usepackage{url}
\usepackage{subfigure}
\usepackage{amsfonts,latexsym} % para tener disponibilidad de diversos simbolos
\usepackage{enumerate}
%\usepackage{booktabs}
\usepackage{float}
%\usepackage{threeparttable}
\usepackage{array,colortbl}
%\usepackage{ifpdf}
%\usepackage{rotating}
\usepackage{cite}
\usepackage{stfloats}
\usepackage{url}
\usepackage{listings}
\usepackage{fancyhdr}

\usepackage{color}
%\usepackage{hyperref}
%\usepackage{wrapfig}
\usepackage{array}
%\usepackage{multirow}
%\usepackage{adjustbox}
%\usepackage{nccmath}
\usepackage[dvipsnames]{xcolor}


\newcommand{\titulo}{Präsenzübungsblatt 7; Statistik}
\newcommand{\nombre}{Lil'Jana}
\newcommand{\FCE}{\textbf{SoSe 2025}}
\newcommand{\dat}{\textbf{Tut 27.05.2025}}

 

\usepackage{fancyhdr}
\pagestyle{fancy}
\fancyhead{}
\fancyfoot{}
\fancyhead[RO]{\small\dat}
\fancyhead[LO]{\small\FCE}
\fancyfoot[RO,LE] {\thepage}

\setcounter{page}{1} %Número de la página, la editorial lo cambiará



\usepackage[onehalfspacing]{setspace}
\begin{document}

\begin{tabular}{p{12cm}}
{{\Large\bf\titulo}}\\
\vspace{0.2cm}\\
 {\large\nombre}\\
\end{tabular}
\vspace{0.2cm}\\
\section{Aufgabe 7.1}
uniformverteilt = gleichverteilt \\
dann: \(
(x,y) : (P((x,y)\in B) = \frac{|B|}{|A|} \text{ mit } |B| = \text{ Fläche von B}
\)\\
Tipp: Dichte von X ausrechnen 
\subsection{a)} gemeinsame Dichte $f_{(x,y)}(x,y) = \frac{\mathbb{1}_A(x,y)}{|A|} = 2 \mathbb{I}_A(x,y)$ \\
Dabei ist $|A|$ die Normierung. 
\[
f_x (x) = \int_{-\infty}^{\infty} f{(x,y)}(x,y) dy = 
\]

Es ist zwar uniformverteilt, aber wir teilen es trotzdem auf, einfach damit wir die Integrale besser lösen können. 
\subsection{b)} 

\subsection{c)} Nein, das geht nicht (man muss ein Gegenbeispiel finden) 

\section{Aufgabe 7.2}

\section{Aufgabe 7.3}

\section{Aufgabe 7.4}

\section{Besprechung}
\section{5.3}
$\alpha \rightarrow$ Erfolgswahrscheinlichkeit 0,05 \\
$\beta \rightarrow$ Erfolgswahrscheinlichkeit 0,6  \\
$\gamma \rightarrow$ 3 Läden mit je Erfolgswahrscheinlichkeit 0,2 (Dabei interessiert uns nicht wie viele Packungen gamma gefunden hat) \\ \\
Gesamt-Erfolgswahrscheinlichkeit für $\gamma: (P\gamma)$ 
\[
\Omega \gamma = \{0,1,2,3\}, \quad \mu \mathcal{B}_{3,(0,2)}, \quad P\gamma = \mathbb{P}(\{1,2,3\}) = 1 - \mathbb{P}(\{0\}) = 1 - \underbrace{\binom{3}{0}}_{1}(0,2)^0 (0,8)^3 = 0,488
\]

\subsection*{a)} nutzen, dass das alles unabhängig ist bei P(w)\\
\[
\Omega = \{(\alpha, \beta, \gamma) \in \{0,1\}^3\} = \{0,1\}^3
\]
Dabei bedeutet 0, dass kein Tofu befunden wurde und 1, dass mind. 1 mal Tofu gefunden wurde. 
\[
\mathbb{P}(\omega) = \prod_{i \in \{\alpha, \beta, \gamma\}} \mu_i (\omega_i) 
\]
mit 
\begin{align*}
    \mu_{\alpha}(1) & = P \alpha = 0,05 & (\mu_{\alpha}(0) = 1-0,05) \\
    \mu_{\beta}(1) & = 0,6 & \vdots \\
    \mu_{\gamma} & = P \gamma = 0,488 & \vdots
\end{align*}

\subsection{b)} hier ist es auch einfacher, wenn man sich das Gegenereignis anschaut, wo keiner Tofu gefunden hat . \\

\end{document}